\section{The cap dilation and its linear limit}\label{sec:dilation}

The cap selects finite solutions whose hot component has a prescribed
weighted size at some time before the horizon. We now magnify that event.
The bath and the quadratic hot interactions disappear under this
magnification; the cold particles converge to two axial velocity atoms
whose densities retain the annular geometry. This is the dilation
argument of \cite[Section~7]{OpenAIPeriodic}, with the additional seed
factor and bath included explicitly below.

Choose horizons \(T_N\downarrow0\), terminal scales
\(s_N/T_N\to0\), and cutoffs \(b_N\to\infty\), and take the cap
solutions furnished by Proposition~\ref{early:cap-selection}. Indices may be
relabeled after passage to a subsequence. Write
\(m_N=m_{T_N}=T_N^p\), and choose a time
\(t_{*,N}\in[0,T_N]\) at which the cap is attained. On the periodic
lift, define
\begin{equation}\label{eq:rescaling}
\begin{gathered}
r_N=t_{*,N}+s_N,\qquad
h_{r_N}(t,x,v)=\frac{h_N(r_Nt,r_Nx,v)}{m_Nr_N},\qquad
C_{r_N}(t,x,v)=r_Nc_N(r_Nt,r_Nx,v),\\
l_N=\frac{s_N}{r_N}\longrightarrow l\in[0,1],\qquad
U_N=\frac{T_N}{r_N}\longrightarrow U\in[1,\infty].
\end{gathered}
\end{equation}
The limits are obtained by subselection, allowing the second ratio to
tend to infinity. We have \(r_N\to0\), \(r_N\le2T_N\) eventually,
and
\(U_N\ge(1+s_N/T_N)^{-1}\), which proves \(U\ge1\).
The rescaled cap and its attained value are
\begin{equation}\label{dil:cap-bound}
0\le h_{r_N}(t,x,v)\le(t+l_N)H(v)\quad(0\le t\le U_N),
\qquad \|h_{r_N}(1-l_N)\|_H=1.
\end{equation}
When no confusion is possible, we suppress the index on \(r_N\).
\begin{figure}[!htbp]
\centering
\resizebox{0.98\linewidth}{!}{%
\begin{tikzpicture}[x=1cm,y=1cm,>=Latex,font=\small,
  timeaxis/.style={->,black!65,line width=.55pt},
  hittick/.style={blue!55!black,line width=1.1pt},
  horizontick/.style={orange!75!black,line width=1.1pt}]
  \path[use as bounding box] (0,-.32) rectangle (13.7,5.97);
  \node[anchor=west,align=left] at (0,5.32)
    {Physical time\\$t$};
  \node[anchor=west,align=left] at (0,3.33)
    {Hit dilation\\$t/r_N$};
  \node[anchor=west,align=left] at (0,1.33)
    {Horizon dilation\\$t/T_N$};

  \draw[black!20,dashed] (2.45,.99)--(2.45,5.37);
  \draw[blue!35!black,dashed] (6.70,.99)--(6.70,5.37);
  \draw[orange!55!black,dashed] (12.05,.99)--(12.05,5.37);
  \foreach \yy in {5.30,3.30,1.30}{
    \draw[timeaxis] (2.20,\yy)--(13.15,\yy);
    \draw[black!60] (2.45,\yy-.10)--(2.45,\yy+.10);
    \draw[hittick] (6.70,\yy-.13)--(6.70,\yy+.13);
    \draw[horizontick] (12.05,\yy-.13)--(12.05,\yy+.13);
    \node[below,fill=white,inner sep=2pt] at (2.45,\yy-.14) {$0$};
  }
  \node[above,text=blue!55!black] at (6.70,5.51) {cap hit};
  \node[above,text=orange!75!black] at (12.05,5.51) {horizon};
  \node[below,fill=white,inner sep=2pt] at (6.70,5.12)
    {$t_{*,N}=r_N-s_N$};
  \node[below,fill=white,inner sep=2pt] at (12.05,5.12) {$T_N$};
  \node[below,fill=white,inner sep=2pt] at (6.70,3.12) {$1-l_N$};
  \node[below,fill=white,inner sep=2pt] at (12.05,3.12) {$U_N=T_N/r_N$};
  \node[below,fill=white,inner sep=2pt] at (6.70,1.12)
    {$\dfrac{1-l_N}{U_N}$};
  \node[below,fill=white,inner sep=2pt] at (12.05,1.12) {$1$};

  \node[align=center,fill=white,inner sep=3pt] at (7.25,4.25)
    {$r_N=t_{*,N}+s_N,\qquad l_N=s_N/r_N$};
  \node[align=center,fill=white,inner sep=3pt] at (7.25,2.24)
    {$t/T_N=(t/r_N)/U_N,\qquad x/T_N=(x/r_N)/U_N$};
  \node[align=center] at (7.05,-.12)
    {$U_N\longrightarrow U<\infty
      \quad\Longrightarrow\quad r_N/T_N\longrightarrow 1/U>0$};
\end{tikzpicture}%
}
\caption{The same hit and horizon in the two dilations.
The rows align physical events and are schematic; the hit may occur at either
endpoint, and no ordering of $r_N$ and $T_N$ is assumed.
Sections~\ref{sec:dilation}--\ref{sec:branching} will show that the hit
dilation has a nonzero weak limit on $0<t<U$ and that $U<\infty$.
The latter conclusion makes the change to the horizon dilation nondegenerate:
compact subsets of $0<t<U$ map inside $0<t<1$.
Thus vanishing in the horizon dilation would force vanishing in the hit
dilation, which is the contradiction used to separate the solutions.}
\label{fig:scales}
\end{figure}


\subsection{The limiting mesh and background}

The dilated annular scales are \(\varrho_{N,j}=s_j/r_N\), \(j\leq N\).
Their ratio is \(e^\eta\), their smallest member is \(l_N\), and
their largest member tends to infinity. Subselect their logarithmic
phase modulo \(\eta\). The grids then converge on compact
subintervals of \((0,\infty)\) to a geometric mesh \(\mathcal M\)
unbounded above. If \(l>0\), its members are
\(l,e^\eta l,e^{2\eta}l,\ldots\). If \(l=0\), it is a two-sided
geometric mesh with no smallest member. To verify this description,
choose a grid point in \((e^{-\eta},1]\), pass to its limit, and
multiply it by the fixed factors \(e^{j\eta}\). When the smallest
point stays positive its integer distance in the grid from that chosen
point is bounded, so can also be fixed by subselection. Coincident phase
endpoints give the same mesh.

For each \(\varrho>0\) and \(\sigma\in\{-1,1\}\), let
\(u_{\varrho,\sigma}(t,z)\) be the unique solution of
\[
u=\sigma+\alpha\tanh\frac{\gamma(z-tu)}\varrho.
\]
Thus \(u_{\varrho,\sigma}=U_{\varrho,\sigma}\) in the notation of the cold profiles. The left side minus the right side has derivative in \(u\) equal to
\(1+t\alpha\gamma S(\gamma(z-tu)/\varrho)/\varrho>0\); its values have
opposite signs for sufficiently large positive and negative \(u\).
This proves existence and uniqueness. The implicit function theorem
gives smoothness, and
\[
|u_{\varrho,\sigma}-\sigma|\leq\alpha,
\qquad
\partial_z u_{\varrho,\sigma}
=\frac{\alpha\gamma S(\gamma(z-tu)/\varrho)}
       {\varrho+t\alpha\gamma S(\gamma(z-tu)/\varrho)}\geq0.
\]
Put \(\theta=1/4\), the fraction of each transverse stripe occupied
by one sign. On an open mesh annulus \(\varrho<|y|<e^\eta\varrho\), define the
nonnegative velocity measure
\begin{equation}\label{eq:background}
\begin{split}
B(t,x,\dd p)
&=\sum_{\sigma=\pm1}b_{\varrho,\sigma}(t,z)
                         \delta_{u_{\varrho,\sigma}(t,z)e_3}(\dd p),\\
b_{\varrho,\sigma}(t,z)
&=\frac{\theta K S(\gamma(z-tu_{\varrho,\sigma})/\varrho)}
 {\varrho+t\alpha\gamma S(\gamma(z-tu_{\varrho,\sigma})/\varrho)}.
\end{split}
\end{equation}
Set \(B=0\) in the hole \(|y|<l\) if \(l>0\), and assign arbitrary
values obeying the same bounds on annular interfaces and the transverse
axis. These exceptional positions have zero Lebesgue measure. In
particular, for every \(k\geq0\),
\begin{equation}\label{dil:background-bound}
\int\langle p\rangle^kB(t,x,\dd p)
\leq\frac{C_k}{t+\varrho}\leq\frac{C_k}{t}
\quad(t>0)
\end{equation}
on an active annulus. If \(l>0\), the bound is instead uniformly
\(C_k/(t+l)\) for all \(t\geq0\).

\begin{proposition}\label{prop:blowup-limit}
Every sequence \eqref{eq:rescaling} has a subsequence for which
\(h_{r_N}\rightharpoonup g\) in local spacetime \(L^1\) on
\((0,U)\times\R^3_x\times\R^3_v\), where
\begin{equation}\label{eq:limit-equation}
0\leq g(t,x,v)\leq(t+l)H(v),\qquad
\D g=2P(B,g)-g\nu_B.
\end{equation}
Here \(B\) is \eqref{eq:background}; its gain and loss are defined by
integration against the velocity measure. Both terms on the right are
locally integrable, with a velocity majorant uniform in position on
each compact positive-time interval. Compactly supported velocity
averages of \(h_{r_N}\) converge strongly in local \(L^1_{t,x}\).
The limit has a locally strongly continuous \(L^1\) representative
at positive times and obeys the integrating-factor formula
\eqref{dil:mild} below.

If \(l>0\), pass also to a limit \(\lambda_N\to\lambda_0\) of the
seed parameters. The conclusion and its local continuity extend to
time zero, with
\begin{equation}\label{dil:initial-trace}
g(0,x,v)=\lambda_0 l\,\chi(x/l,v).
\end{equation}
\end{proposition}

\subsection{Cold convergence: survival, Jacobian, and stripe averaging}

We first identify the cold measure independently of the hot limit.
Fix a compact time-space region
\[
a\leq t\leq A<U,\qquad |x|\leq R,
\]
where \(a>0\). Since \(U_N\to U\) and \(A<U\), eventually
\(r_NA<T_N\). The cap therefore holds on the entire initial
physical interval \([0,r_NA]\), as required for the survival
estimate below.
By Proposition~\ref{cold:control}, the dilated
cold velocity moments are bounded there, and indeed uniformly in
position, by
\begin{equation}\label{dil:cold-moment-bound}
\int\langle p\rangle^k C_{r_N}(t,x,p)\dd p
\leq C_k\left(\frac1{t+l_N}+1\right).
\end{equation}
The nonmain tails have the still smaller bound
\(C_M r_N^{M+1}t^Me^{-4|p|^2}\) for any fixed \(M>0\).
Uniformly decreasing velocity tails hold under these moment bounds.

First restrict to annuli with \(\varrho_{N,j}\) in a fixed compact
interval \([\varepsilon,R_1]\subset(0,\infty)\). Their number is
bounded independently of \(N\). A free-core path at scale
\(s=r_N\varrho_{N,j}\) has bounded velocity and remains in its own
spatial enlargement. At its position, all free velocities in that
color are within \(C d_s\) of the same center, so the contribution
to its loss frequency is at most \(C d_s/s\).
The remainder in Proposition~\ref{cold:control} has velocity moment
amplitude \(O(s^{15/2})\), by Lemma~\ref{cold:core-amplitude}.
The nonmain colors have the bound of Lemma~\ref{cold:nonmain}; the
bath bound and the cap bound give bounded frequencies at these
velocities. Hence, over physical times \(0\le t_{\rm old}\le r_NA\),
the optical depth on every free-core path is at most
\[
 C_A r_N\bigl(d_s/s+s^{15/2}+1\bigr)=o(1)
\]
uniformly on this compact scale range. The cutoff only decreases
this frequency. Keeping the freely transported term in the positive
path formula of \eqref{con:colors} gives
\[
 e^{-\epsilon_N}f_i\le c_i\le f_i+\epsilon_s m_i,
 \qquad \epsilon_N\longrightarrow0,
\]
where \(f_i\) is the free core of \eqref{cold:free-core}; the symbols
\(\epsilon_N\) and \(\epsilon_s\) denote different quantities.
After multiplication by \(r_N\), the remainder has velocity mass
\begin{equation}\label{dil:remainder-power}
 r_N O(s^{15/2})=O(r_N^{17/2}\varrho_{N,j}^{15/2}).
\end{equation}
The same estimates hold with every fixed polynomial velocity weight.
Consequently the dilated actual color and its free core differ by
\(o(1)\) in velocity total variation, uniformly on the compact
regions under consideration.

On bounded dilated spatial sets, only the copy centered at the origin
can be main for sufficiently large \(N\). Indeed, an intersecting
copy would have a lattice center \(n\) with
\(|\bar n|\le s_0+o(1)\) and
\(|n_3|\le 1/8+b_{\rm sp}+o(1)\). Since \(1/8+b_{\rm sp}=3/16\)
and \(s_0\) is small, this contradicts \(|n|\ge1\) unless
\(n=0\). This argument uses the Euclidean spacing of the rotated
lattice and identifies the free cores directly on its lift.

We compute a centered free core exactly. Fix one sign and one such annulus,
write \(\varrho_N=s/r_N\), and use the velocity parameters
\[
\bar p=d_s\bar w,\qquad
p_3=\sigma+\alpha\tanh\frac{\gamma(z-tp_3)}{\varrho_N}+d_sw_3.
\]
The axial map depends only on \(w_3\). Its derivative is
\[
\frac{\partial p_3}{\partial w_3}
=\frac{d_s}
 {1+t\alpha\gamma S(\gamma(z-tp_3)/\varrho_N)/\varrho_N}.
\]
Thus Equation~\eqref{cold:core-formula}, after dilation, says that the
velocity integral of the free core against a test \(\psi\) is
\begin{equation}\label{dil:free-core-integral}
\begin{split}
\int r_N f_i(r_Nt,r_Nx,p)\psi(p)\dd p
={}&\int J(w)
 \frac{K S(\gamma(z-tp_3)/\varrho_N)}
      {\varrho_N+t\alpha\gamma S(\gamma(z-tp_3)/\varrho_N)}
 \psi(d_s\bar w,p_3)\\
&\quad\cdot\ind_{E_j^\sigma}(r_Ny-r_Ntd_s\bar w)
 \ind_{\{|r_N(z-tp_3)|<1/8\}}\dd w.
\end{split}
\end{equation}
On the compact regions under consideration the last axial indicator
equals one for all sufficiently large \(N\). Since the implicit
velocity map has derivative at least one in \(p_3\),
\[
|p_3-u_{\varrho_N,\sigma}(t,z)|\leq d_s|w_3|.
\]
After the mesh subselection, \(\varrho_N\to\varrho>0\), so the velocity
test and the smooth amplitude in \eqref{dil:free-core-integral}
converge uniformly for \(w\in\supp J\) to
\[
\psi(u_{\varrho,\sigma}e_3)
\frac{K S(\gamma(z-tu_{\varrho,\sigma})/\varrho)}
 {\varrho+t\alpha\gamma S(\gamma(z-tu_{\varrho,\sigma})/\varrho)}.
\]

It remains to average the transverse indicator. In dilated coordinates
its period is
\[
\frac{s^2}{r_N}=r_N\varrho_N^2\longrightarrow0.
\]
The transverse shift of its argument is \(td_s\bar w\), whose ratio
to that period is at most
\[
\frac{Ad_s}{r_N\varrho_N^2}=A d_0 r_N^9\varrho_N^8\longrightarrow0.
\]
For the indicator of the fixed periodic stripe pattern, a translation
by this vanishing relative amount changes its integral on a bounded
set by \(o(1)\), uniformly in its phase and in bounded \(t,w\).
One can verify this by summing the lengths of the small intervals
around stripe endpoints. The annular margins have dilated width
\(O(s^2/r_N)=O(r_N\varrho_N^2)\), so their area also tends to zero on
the annulus. The outer annular radii converge; their boundaries can
likewise be removed in sets of arbitrarily small area.

Finally the periodic stripe indicators converge weak-* to their mean
\(\theta\). Explicitly, if \(q\) is the one-period indicator and
\(q-\theta\) has bounded periodic primitive \(Q\), then for a smooth
compactly supported one-dimensional test \(f\), period \(p\), and
arbitrary phase \(\zeta\),
\[
\left|\int f(y_1)(q(y_1/p+\zeta)-\theta)\dd y_1\right|
\leq p\|Q\|_\infty\|f'\|_1.
\]
Fubini's theorem gives the same assertion with the other coordinates
and parameters present. Approximation in \(L^1\) extends it to the
annular restrictions. Integrating \(J\), whose mass is one, in
\eqref{dil:free-core-integral} produces exactly the coefficient
\(b_{\varrho,\sigma}\) in \eqref{eq:background}.

These arguments identify the cold limit on all compact positive
scale annuli. To remove the lower scale restriction, use
\eqref{dil:cold-moment-bound}: on the slab \(t\geq a\), the integral
over \(|y|<C\varepsilon\), with bounded \(z\), is at most
\(C_{a,A,R,k}\varepsilon^2\). The same estimate holds for the
candidate \(B\). Scales larger than a fixed multiple of \(R\) have
no main color in the chosen spatial compact set, and their escaped
tails vanish. We have proved that for every continuous compactly
supported velocity test \(\psi\),
\begin{equation}\label{dil:cold-weakstar}
\int C_{r_N}(t,x,p)\psi(p)\dd p
\stackrel{*}{\rightharpoonup}
\int\psi(p)B(t,x,\dd p)
\quad\hbox{in local }L^\infty_{t,x}
\end{equation}
on \((0,U)\times\R^3\). Indeed the computation first gives weak
convergence against smooth compact time-space tests, and the common
\(L^\infty\) bound extends it to all \(L^1\) tests. Polynomial
velocity tails can also be removed uniformly. If \(l>0\), choose
\(\varepsilon<l/2\); all the bounds used above hold down to time
zero, so \eqref{dil:cold-weakstar} does too.

\subsection{The hot equation and the size of every error}

The transport derivative of \(h_N(r_Nt,r_Nx,v)/(m_Nr_N)\)
is \(m_N^{-1}(\D h_N)(r_Nt,r_Nx,v)\). Thus the equation for
\(h_N\) gives exactly
\begin{equation}\label{dil:prelimit-equation}
\D h_{r_N}
=2P^{b_N}(C_{r_N},h_{r_N})
 -h_{r_N}\nu^{b_N}_{C_{r_N}}+E_N,
\end{equation}
where, writing \(u_N^r(t,x,v)=u_N(r_Nt,r_Nx,v)\),
\[
\begin{split}
E_N={}&r_N\bigl[2P^{b_N}(u_N^r,h_{r_N})
                   -h_{r_N}\nu^{b_N}_{u_N^r}\bigr]\\
 &+m_Nr_N^2\bigl[P^{b_N}(h_{r_N},h_{r_N})
                   -h_{r_N}\nu^{b_N}_{h_{r_N}}\bigr]
   +m_N^{-1}\mathcal S_N(r_Nt,r_Nx,v).
\end{split}
\]
Here \(\mathcal S_N\) is the nonnegative source in
\eqref{con:colors}, including the cold--bath and distinct-color gains.
The bath estimate \eqref{early:bath-bound} gives \(u_N^r\le C H\). For every fixed
\(A<\infty\), the cap and the gain estimates therefore imply, on
\(0\le t\le\min(A,U_N)\),
\begin{equation}\label{dil:error-bounds}
 |E_N(t,x,v)|
 \le C_A\bigl(r_N+m_Nr_N^2+T_N^{M-1-p}\bigr)
             \langle v\rangle H(v),
 \qquad M>p+1.
\end{equation}
For the last term we used the arbitrary-power source bound of
\eqref{early:source}:
\[
 \frac{\mathcal S_N(r_Nt,r_Nx,v)}{m_N}
 \le C_M\frac{r_N^{M-1}}{T_N^p}t^{M-1}H(v)
 \le C_{A,M}T_N^{M-1-p}H(v).
\]
Thus every error tends to zero locally in spacetime \(L^1\), including
down to time zero on bounded dilated intervals. The use of
\(r_N\le2T_N\) here requires no positive lower bound for
\(r_N/T_N\).

The main cold sheets and Lemma~\ref{lem:sheet-bound} also give
\[
P^{b_N}(C_{r_N},H)\leq
 C\left(\frac1{t+l_N}+1\right)H,
\qquad
\nu^{b_N}_{C_{r_N}}(v)\leq
 C\left(\frac1{t+l_N}+1\right)\langle v\rangle.
\]
The nonmain cold component is bounded by a vanishing multiple of
\(e^{-4|v|^2}\leq CH\), so its gain is covered by
Lemma~\ref{lem:hot-bound}. Equations \eqref{dil:cap-bound} and
\eqref{dil:prelimit-equation} now imply, on \(a\leq t\leq A<U\),
\begin{equation}\label{dil:transport-majorant}
|\D h_{r_N}(t,x,v)|\leq C_{a,A}\langle v\rangle H(v).
\end{equation}
This is a genuine integrable velocity majorant, uniform in position.
If \(l>0\), the same conclusion holds on \(0\leq t\leq A\), since
\(l_N\geq l/2\) eventually.

The cap gives weak compactness of \(h_{r_N}\) in local \(L^1\),
and \eqref{dil:transport-majorant} permits application of
Lemma~\ref{comp:averaging}. Subselect on an exhaustion of
\((0,U)\times\R^3\). The resulting limit \(g\) satisfies the
bound in \eqref{eq:limit-equation}, and each compactly supported
velocity average converges strongly in local \(L^1_{t,x}\).

\subsection{Passing the cold--hot products}

The dilated cold densities concentrate in velocity, so the
two-density compactness argument of Section~\ref{sec:compactness}
does not apply to \(C_{r_N}\). The needed product passage uses
\eqref{dil:cold-weakstar} and hot averaging instead.

For compact velocity tests \(\psi,\chi\), put
\[
a_N(t,x)=\int\psi(p)C_{r_N}(t,x,p)\dd p,\qquad
\zeta_N(t,x)=\int\chi(v)h_{r_N}(t,x,v)\dd v.
\]
On every compact positive-time region, \(a_N\) is bounded in
\(L^\infty\) and converges weak-* to
\(a=\int\psi\dd B\), while
\(\zeta_N\to\zeta=\int\chi g\dd v\)
strongly in \(L^1\). For a bounded time-space test \(\varphi\),
\[
\int\varphi(a_N\zeta_N-a\zeta)
=\int\varphi a_N(\zeta_N-\zeta)
 +\int\varphi(a_N-a)\zeta\longrightarrow0.
\]
Finite sums of separated velocity tests approximate continuous tests
uniformly on compact velocity products. The total product mass is
uniformly bounded on compact time-space regions by
\eqref{dil:cold-moment-bound} and the cap. We conclude
\begin{equation}\label{dil:mixed-product}
C_{r_N}(t,x,p)h_{r_N}(t,x,v)\dd p\dd v\dd x\dd t
\rightharpoonup B(t,x,\dd p)g(t,x,v)\dd v\dd x\dd t
\end{equation}
against continuous compact velocity tests, with arbitrary smooth
compact time-space dependence. Uniform moment bounds in both velocities
allow removal of velocity cutoffs in tests with any fixed polynomial
growth for which they are used below.

Here is an explicit weak collision test suitable for this convergence.
For a continuous output test \(\phi\), let
\[
\mathcal K_\phi(p,v)
=\frac{|p-v|}{2}\int_{S^2}
 \phi\left(\frac{p+v}{2}+\frac{|p-v|}{2}\omega\right)\dd\omega.
\]
Lemma~\ref{lem:carleman} gives
\(\int\phi P(c,h)=\iint c(p)h(v)\mathcal K_\phi(p,v)\dd p\dd v\).
The test \(\mathcal K_\phi\) is continuous in both incoming
velocities, including \(p=v\), and is bounded in absolute value by
\(2\pi|p-v|\|\phi\|_\infty\). The loss test is
\(2\pi|p-v|\phi(v)\), with the same properties. Inserting smooth
time-space dependence in \(\phi\) does not change these facts.
On bounded velocity products the factor
\(\min(1,b_N/|p-v|)\) is eventually one; its complement is
uniformly negligible by the weighted product tails. Applying
\eqref{dil:mixed-product} proves the required gain and loss limits.

The limiting gain is an integrable function. To check the possible
exception at the velocity axis, use the sphere formula: every
nondegenerate diameter sphere meets that axis in at most two points,
which have zero sphere-area measure. A degenerate sphere has zero
collision rate. Tonelli's theorem therefore gives no output mass on
the axis. Off the axis, the collision-plane formula of
Lemma~\ref{lem:carleman} gives a density for each axial partner.
Each summand of \(B\) is an allowed zero-width axial sheet in
Lemma~\ref{lem:sheet-bound}; hence on
\(a\leq t\leq A\),
\[
0\leq P(B,g)\leq C_{a,A}H,
\qquad 0\leq g\nu_B\leq C_{a,A}\langle v\rangle H.
\]
These bounds dominate the entire weak gain measure, including across
the velocity axis. Passing \eqref{dil:prelimit-equation}, using
\eqref{dil:error-bounds}, proves the equation in
\eqref{eq:limit-equation}.

\subsection{A positive lower mesh endpoint and the local mild formula}

When \(l>0\), \eqref{dil:transport-majorant} holds uniformly down to
zero. The rescaled initial hot density, on every fixed spatial compact
set for large \(N\), is exactly
\[
h_{r_N}(0,x,v)=\lambda_N l_N\chi(x/l_N,v).
\]
The other periodic copies have moved outside that compact set.
Subselect \(\lambda_N\to\lambda_0\). Smoothness and compact support
of \(\chi\), and \(l_N\to l>0\), give strong local \(L^1\)
convergence to \eqref{dil:initial-trace}. Passing the weak transport
identity with tests touching time zero is justified by the uniform
source majorant and yields that initial trace for \(g\). The cold
product argument also extends to zero in this case, as already noted.
For \(l=0\) the coefficient can be of order \(1/t\), and the
proposition only asserts the equation on the open positive-time
interval.

On any \([a,A]\subset(0,U)\), the limiting transport source is
locally integrable and has the preceding velocity majorant uniformly
in position. Changing to free coordinates, first on bounded velocities and enlarged
spatial compact sets, gives a locally strongly continuous \(L^1\)
representative. The uniform integrable velocity tails then remove the
velocity truncation. The cap bound passes to every trace of this
representative. Almost every free path obeys the scalar absolutely
continuous equation
\(\D g+\nu_Bg=2P(B,g)\). Its loss coefficient is locally integrable
along the path, since
\(\nu_B(t,x,v)\leq C_a\langle v\rangle\) for \(t\geq a\).
The integrating factor therefore gives, for \(0<a<t<U\),
\begin{equation}\label{dil:mild}
\begin{split}
g(t,x,v)={}&g(a,x-(t-a)v,v)
 \exp\left[-\int_a^t\nu_B(\tau,x-(t-\tau)v,v)\dd\tau\right]\\
&+2\int_a^t P(B,g)(s,x-(t-s)v,v)
 \exp\left[-\int_s^t\nu_B(\tau,x-(t-\tau)v,v)\dd\tau\right]\dd s.
\end{split}
\end{equation}
This is an almost-everywhere identity, or equivalently an identity in
local \(L^1\) at the displayed times. It holds with \(a=0\) when
\(l>0\). Altering the medium on the null annular interfaces does not
change the identity: free transport preserves phase-space measure, so
Fubini's theorem discards the affected set of paths. These statements
complete the proof of Proposition~\ref{prop:blowup-limit}.

\subsection{The free-core comparison used for nonconcentration}

Weak convergence alone need not preserve the attained weighted norm in
\eqref{dil:cap-bound}. The next section will exclude its concentration
near the velocity axis or in a small set of collision histories. Its
input is an upper comparison involving only the explicit free cores.

\begin{lemma}\label{dil:free-core-comparison}
For all sufficiently large \(N\), define the velocity measure
\[
 A_N(t,x,\dd p)
   =r_N\sum_{i:\,j\le N}f_i(r_Nt,r_Nx,p)\dd p,
       \qquad 0\le t\le1.
\]
This definition uses free transport of the initial colors and is
meaningful even if \(t>U_N\). At each position, at most one summand
is nonzero. Its velocity moments and normalized gain obey
\[
 \int\langle p\rangle^k A_N(t,x,\dd p)
      \le\frac{C_k}{t+l_N},\qquad
 P(A_N,H)\le\frac{C}{t+l_N}H.
\]
There are numbers \(\varepsilon_N\downarrow0\) such that, on
\(0\le t\le\min(1,U_N)\), the hot component obeys the pathwise
upper comparison
\begin{equation}\label{dil:positive-comparison}
 \D h_{r_N}\le2P(A_N,h_{r_N})+\varepsilon_NH.
\end{equation}
The bounds and comparison hold on full-measure endpoint sets at each
time section, with representatives given by their path formulas.
\end{lemma}

\begin{proof}
The supports of the free cores lie in their disjoint spatial
enlargements through physical time \(r_N\); no retirement occurs
on this interval for large \(N\). The exact Jacobian formula
\eqref{dil:free-core-integral}, or its version before centering a
periodic copy, gives the moment estimate because
\(r_N/(r_Nt+s_j)\le1/(t+l_N)\). The sheet gain estimate gives
the second bound.

In the upper cold estimate, the main-color remainder contributes at
most \(C r_Ns_j^{15/2}H\) to its gain against \(H\).
The sum of the nonmain contributions has a vanishing Gaussian
majorant after multiplication by \(r_N\). Uniformly over all
positions, therefore,
\[
 P(C_{r_N},H)\le P(A_N,H)+o(1)H.
\]
More precisely this inequality follows by applying the positive gain
to the nonnegative upper error in \(C_{r_N}\le A_N+R_N\), where
\(P(R_N,H)\le o(1)H\); it does not assert a signed total-variation
comparison at all scales. On the present time range
\(h_{r_N}\le2H\). Drop every nonnegative loss in
\eqref{dil:prelimit-equation}, use \(P^{b_N}\le P\), and bound the
positive bath, hot, and source terms exactly as in
\eqref{dil:error-bounds}. This proves
\eqref{dil:positive-comparison} with an error independent of
\(t,x,v\).

All finite approximations obey their equations in free coordinates.
Fubini's theorem gives the comparisons along almost every free path
and on full-measure endpoint sets at every specified terminal time.
The countably many approximation indices and lattice translations
can be treated simultaneously. This is the form of the comparison
used in the next section.
\end{proof}
